% Propietario conservado: /Users/ruben/Documents/New project/output/EXTREMA_MEDIA_RAZON_RECIPROCIDAD_ES_EN_20260918/ARTICULO/ES/source/libro/03_refinamiento.tex
\chapter{Raffinement conservatif de l’unité et mémoire des retenues}
\label{x_reciprocidad:lib03:refinamiento}

La complétion dimensionnelle et le registre des opérations admettent une composition arithmétique explicite. L'architecture APP--TRIT--TPK fournit préalablement l'état enrichi, les lectures orientées et les complétions de l'unité. Sur ces sorties, on construit ici un transport affine qui conserve la somme normalisée et permet de récupérer, par divisions euclidiennes, tout mot fini de retenues canoniques. Sa limite possède des cylindres de diamètre contrôlé ; les marques de frontière distinguent les histoires qu'une évaluation scalaire identifie.

Le bloc central, ses concaténations et la partition
\(1000=729+243+27+1\) sont des antécédents de cette construction. La composition
de ces objets avec le transport affine, l’inverse par les restes et son
prolongement constitue la formalisation développée dans cette section.
L’application agit sur un domaine arithmétique déclaré. Chaque réalisation
sur une histoire enrichie particulière conserve en outre la règle
prospective qui produit ses retenues.

\section{Relation entre norme quadratique et complétion cubique}
\label{x_reciprocidad:lib03:completacion}

Soient \(b\geq1\) entier et \(B=b+1\geq2\). Nous définissons
\begin{equation}
 p_b=b^2-b+1,\qquad q_b=B^2-p_b=3b.
 \label{x_reciprocidad:lib03:pb}
\end{equation}

\begin{proposition}[Transfert de la complétion]
\label{x_reciprocidad:lib03:prop-completacion}
Les identités
\begin{equation}
 Bp_b=b^3+1,\qquad Bq_b+1=B^3-b^3
 \label{x_reciprocidad:lib03:cuadratica-cubica}
\end{equation}
déterminent la transformation
\begin{equation}
 (p_b,q_b)\longmapsto(Bp_b-1,Bq_b+1)
       =(b^3,B^3-b^3).
 \label{x_reciprocidad:lib03:transferencia-cubica}
\end{equation}
La somme passe de \(B^2\) à \(B^3\), tandis que la somme des
coordonnées divisées par leurs échelles respectives reste égale à un.
\end{proposition}

\begin{proof}
La factorisation \(b^3+1=(b+1)(b^2-b+1)\) fournit la première
identité. La seconde résulte de
\((b+1)^3-b^3=3b^2+3b+1\). Les additionner prouve le changement d’échelle et,
après division par \(B^3\), la conservation de l’unité normalisée.
\end{proof}

Para \(b=9\),
\begin{equation}
 (73,27)\longmapsto(10\cdot73-1,10\cdot27+1)=(729,271).
 \label{x_reciprocidad:lib03:caso-729}
\end{equation}
Le complément conserve la décomposition \(271=243+27+1\). Dans la complétion précédente, la dernière contribution correspond à l'apex ajouté. La correction \((-1,+1)\) a une somme nulle, et la multiplication par dix transforme l'échelle de cent en celle de mille. Pour la famille complète, le transport des proportions est
\begin{equation}
 \left(\frac{p_b}{B^2},\frac{q_b}{B^2}\right)
 \longmapsto
 \left(\frac{p_b}{B^2}-\frac1{B^3},
       \frac{q_b}{B^2}+\frac1{B^3}\right).
 \label{x_reciprocidad:lib03:fracciones-completacion}
\end{equation}
La quantité conservée dans cette identité est une somme linéaire. Les formes
quadratiques et les grandeurs physiques possèdent leurs lecteurs propres, qui
seront composés ultérieurement avec ce transport.

\section{Concaténations du bloc central et restitution unitaire}
\label{x_reciprocidad:lib03:centro}

Pour \(b\geq2\), les chiffres \(b-2\) et \(2\) appartiennent à l’alphabet de
base \(B=b+1\). Considérons la famille algébrique
\begin{equation}
 B_c^{(b)}=\begin{pmatrix}b-2&2\\2&b-2\end{pmatrix},\qquad
 C_B(u,v)=Bu+v.
 \label{x_reciprocidad:lib03:bloque-central}
\end{equation}
Le cas \(b=9\) retrouve le bloc central documenté
\(\bigl(\begin{smallmatrix}7&2\\2&7\end{smallmatrix}\bigr)\), de
valeurs propres \(9\) et \(5\). Leurs concaténations vérifient
\(72+27=77+22=99\). La concaténation, le spectre et la position centrale sont des
lectures distinctes du même antécédent et maintiennent leurs domaines.

L’application du lecteur aux lignes de \eqref{x_reciprocidad:lib03:bloque-central} donne
\begin{align}
 C_B(b-2,2)&=B(b-2)+2=b(b-1)=p_b-1,\\
 C_B(2,b-2)&=2B+(b-2)=3b=q_b.
 \label{x_reciprocidad:lib03:concatenaciones}
\end{align}
La somme est \(B^2-1\). La restitution d'une unité dans la première coordonnée et le transfert suivant produisent
\begin{equation}
 (p_b-1,3b)\xrightarrow{\ +(1,0)\ }(p_b,3b)
 \xrightarrow{\ L_1\ }(b^3,B^3-b^3).
 \label{x_reciprocidad:lib03:cadena-restauracion}
\end{equation}
En particulier,
\begin{equation}
 (72,27)\longmapsto(73,27)\longmapsto(729,271).
 \label{x_reciprocidad:lib03:cadena-central}
\end{equation}
La première flèche augmente la somme de un. La seconde change l’échelle
et redistribue une unité avec une correction de somme nulle. Pour représenter
explicitement la contribution avant son incorporation, le triplet
\((p_b-1,3b,1)\) a déjà pour somme \(B^2\), et l’application
\((x,y,r)\mapsto(x+r,y,0)\) conserve cette somme. Son inversion requiert de
conserver la donnée \(r=1\) : sur des triplets arbitraires, l’application n’est
pas injective.

La restitution est ainsi enregistrée sans lui attribuer une conservation fermée de la paire avant l'incorporation de l'unité. La famille en \(b\) étend algébriquement le cas central reçu ; la sélection des états de l'atlas reste celle établie par APP--TRIT--TPK. La composition \eqref{x_reciprocidad:lib03:cadena-restauracion} détermine une retenue égale à un dans cette transition, en laissant ouverte la spécification prospective des opérations suivantes.

\section{Transport affine et domaine de positivité}
\label{x_reciprocidad:lib03:afin}

\begin{definition}[Transport avec retenue]
\label{x_reciprocidad:lib03:def-transporte}
Pour une base entière \(B\geq2\) et \(c\in\mathbb Z\), on définit
\begin{equation}
 L_c(x,y)=(Bx-c,By+c).
 \label{x_reciprocidad:lib03:lc}
\end{equation}
L’entier \(c\) spécifie le transfert orienté entre les deux
coordonnées. Son signe positif réduit la première par rapport à \(Bx\)
et augmente la seconde de la même quantité.
\end{definition}

\begin{proposition}[Inverse et positivité]
\label{x_reciprocidad:lib03:prop-inversa}
Dans la réalisation affine ultérieure \(\mathbb R^2\), avec \(B,c\) fixés,
\begin{equation}
 L_c^{-1}(X,Y)=\left(\frac{X+c}{B},\frac{Y-c}{B}\right),\qquad
 X+Y=B(x+y).
 \label{x_reciprocidad:lib03:inversa-afin}
\end{equation}
L’antécédent est entier exactement lorsque \(X+c\) et \(Y-c\) sont
divisibles par \(B\). Pour \(x,y\geq0\), l’image est non négative
exactement lorsque
\begin{equation}
 -By\leq c\leq Bx.
 \label{x_reciprocidad:lib03:positividad}
\end{equation}
\end{proposition}

\begin{proof}
La substitution directe démontre les deux identités de
\eqref{x_reciprocidad:lib03:inversa-afin} et son critère d’intégralité. Les
inégalités \(Bx-c\geq0\), \(By+c\geq0\) équivalent à
\eqref{x_reciprocidad:lib03:positividad}.
\end{proof}

Pour l'alphabet canonique \(c\in\{0,\ldots,B-1\}\), tous les prolongements sont admissibles si \(x\geq1\), \(y\geq0\). À la frontière \(x=0\), la non-négativité n'admet que \(c=0\). Le domaine positif et le domaine entier général sont ainsi différenciés par une condition explicite.

\Needspace{12\baselineskip}
\section{Récupération canonique au moyen du résidu}
\label{x_reciprocidad:lib03:residuo}

\begin{theorem}[Inverse canonique par résidu]
\label{x_reciprocidad:lib03:teo-residuo}
Soient \(M\geq1\) entier et \(X,Y\geq0\) entiers tels que
\(X+Y=BM\). Il existe un unique triplet admissible \((x,y,c)\) qui satisfait
\[
 x+y=M,\qquad 0\leq c\leq B-1,\qquad L_c(x,y)=(X,Y).
\]
Il s’obtient par
\begin{equation}
 c=Y\bmod B=(-X)\bmod B,\qquad
 y=\frac{Y-c}{B},\qquad x=\frac{X+c}{B}.
 \label{x_reciprocidad:lib03:residuo-inversor}
\end{equation}
\end{theorem}

\begin{proof}
Tout antécédent vérifie \(Y\equiv c\pmod B\). Le représentant dans
\(\{0,\ldots,B-1\}\) est unique. La congruence \(X+Y\equiv0\pmod B\)
implique la divisibilité de \(X+c\). La division euclidienne de \(Y\)
produit \(y\geq0\), et \(X+c\geq0\) produit \(x\geq0\). Si
\(c>0\), le cas \(x=0\) exigerait \(X=-c<0\), de sorte que le
triplet obtenu appartient au domaine admissible. Sa somme est \(M\), et
la substitution dans \eqref{x_reciprocidad:lib03:lc} retrouve \((X,Y)\). L’unicité
du résidu et des deux quotients prouve l’unicité du triplet.
\end{proof}

Le couple entier et son échelle contiennent la dernière retenue canonique. La
récupération dérive de son reste et ne requiert pas de conserver une autre copie
indépendante de cette retenue. Il y a \(M+1\) triplets avec \(c=0\) et
\(M\) triplets dans chacune des \(B-1\) classes restantes : le total
est \(BM+1\), égal au nombre de couples non négatifs de somme \(BM\).

Pour une retenue entière générale, écrivons \(c=r+B\ell\), avec \(0\leq r\leq B-1\). Le résidu détermine \(r\), tandis que la reconstruction complète utilise également l'élévation \(\ell\) :
\begin{equation}
 y=\frac{Y-r}{B}-\ell,\qquad
 x=\frac{X+r}{B}+\ell.
 \label{x_reciprocidad:lib03:elevacion}
\end{equation}
La conservation de l'élévation distingue la retenue signée de sa réduction résiduelle.

\section{Itération variable et mémoire finie complète}
\label{x_reciprocidad:lib03:memoria}

Fixons \((x_0,y_0)\in\mathbb Z^2\), de somme \(M_0>0\), et une suite d'opérations \(c_1,c_2,\ldots\). Nous définissons
\begin{equation}
 (x_n,y_n)=L_{c_n}\cdots L_{c_1}(x_0,y_0),\qquad
 C_n=\sum_{j=1}^n c_jB^{n-j},\qquad C_0=0.
 \label{x_reciprocidad:lib03:registro-acumulado}
\end{equation}
Les retenues sont reçues de la règle d'opération correspondante ; la valeur limite n'intervient pas dans leur sélection.

\begin{theorem}[Registre accumulé exact]
\label{x_reciprocidad:lib03:teo-acumulacion}
Pour tout \(n\geq0\),
\begin{equation}
 \begin{aligned}
 x_n&=B^nx_0-C_n,& y_n&=B^ny_0+C_n,\\
 x_n+y_n&=B^nM_0,& C_{n+1}&=BC_n+c_{n+1}.
 \end{aligned}
 \label{x_reciprocidad:lib03:iteracion-exacta}
\end{equation}
\end{theorem}

\begin{proof}
Pour \(n=0\), les identités sont immédiates. Si elles sont valables en
\(n\), l’application de \(L_{c_{n+1}}\) multiplie par \(B\) toutes
les contributions précédentes et ajoute \(c_{n+1}\) au cumul.
On obtient la récurrence et, par substitution, les deux coordonnées de
l’étape suivante. La somme élimine \(C_n\).
\end{proof}

Lorsque les retenues sont canoniques,
\begin{equation}
 0\leq C_n\leq B^n-1.
 \label{x_reciprocidad:lib03:intervalo-acumulado}
\end{equation}
En particulier, \(x_0\geq1\), \(y_0\geq0\) garantissent
\(x_n\geq B^n(x_0-1)+1\) et \(y_n\geq0\) à toutes les étapes.

\begin{corollary}[Récupération du mot complet]
\label{x_reciprocidad:lib03:cor-palabra}
Étant connus \(B,n\) et le couple final, une histoire canonique finie
récupère son entrée et toutes ses retenues au moyen de
\begin{align}
 C_n&=y_n\bmod B^n,&
 y_0&=\left\lfloor\frac{y_n}{B^n}\right\rfloor,&
 x_0&=\frac{x_n+C_n}{B^n},
 \label{x_reciprocidad:lib03:recuperacion-inicial}\\
 c_j&=\left\lfloor\frac{C_n}{B^{n-j}}\right\rfloor\bmod B,
 &&1\leq j\leq n.
 \label{x_reciprocidad:lib03:recuperacion-digitos}
\end{align}
\end{corollary}

\begin{proof}
La borne \eqref{x_reciprocidad:lib03:intervalo-acumulado} fait de \(y_n=B^ny_0+C_n\) une division euclidienne. Son quotient et son reste fournissent \(y_0,C_n\) ; la première coordonnée donne \(x_0\). Les divisions successives de \(C_n\) par des puissances de \(B\) récupèrent les chiffres. De manière équivalente, on applique \eqref{x_reciprocidad:lib03:residuo-inversor} depuis la dernière étape jusqu'à la première.
\end{proof}

La profondeur conserve les zéros initiaux du mot, même s’ils
n’altèrent pas l’entier \(C_n\). Si l’on conserve \(M_0\) au lieu de
\(n\), l’égalité exacte \((x_n+y_n)/M_0=B^n\) récupère la
profondeur. Normaliser le couple pour que sa somme soit un élimine ce compteur
d’échelle lorsqu’il n’est pas conservé par une autre lecture.

\section{Composition de blocs et compatibilité avec la partition}
\label{x_reciprocidad:lib03:composicion}

Un bloc de \(n\) opérations, de cumul \(C\), agit par
\begin{equation}
 L_C^{[n]}(x,y)=(B^nx-C,B^ny+C).
 \label{x_reciprocidad:lib03:bloque}
\end{equation}
La composition de deux blocs vérifie
\begin{equation}
 L_D^{[m]}\circ L_C^{[n]}=L_{B^mC+D}^{[n+m]}.
 \label{x_reciprocidad:lib03:ley-bloques}
\end{equation}
En effet, substituer la première image dans la seconde multiplie son cumul par \(B^m\) et ajoute \(D\). La loi
\[
 (n,C),(m,D)\longmapsto(n+m,B^mC+D)
\]
est associative par composition d’applications. Elle exprime la contribution
du premier bloc à l’échelle du second et maintient la sortie lors du
regroupement des mêmes événements ordonnés. Ajouter de nouveaux événements modifie
le cumul selon cette même loi.

\section{Convergence et contrôle quantitatif de la lecture}
\label{x_reciprocidad:lib03:limite}

Écrivons \(M_n=B^nM_0\) et
\begin{equation}
 f_n=\frac{x_n}{M_n},\quad g_n=\frac{y_n}{M_n},\quad
 \theta_n=\frac{C_n}{B^n}=\sum_{j=1}^n\frac{c_j}{B^j}.
 \label{x_reciprocidad:lib03:lecturas-normalizadas}
\end{equation}
Les identités exactes sont
\begin{equation}
 f_n=\frac{x_0}{M_0}-\frac{\theta_n}{M_0},\qquad
 g_n=\frac{y_0}{M_0}+\frac{\theta_n}{M_0},\qquad f_n+g_n=1.
 \label{x_reciprocidad:lib03:unidad-normalizada}
\end{equation}

\begin{theorem}[Convergence des retenues bornées]
\label{x_reciprocidad:lib03:teo-limite}
Si \(|c_j|\leq C\) pour tout \(j\), la série \(\theta=\sum_{j\geq1}c_jB^{-j}\) converge absolument. Les lectures \(f_n,g_n\) convergent vers \(f,g\), avec \(f+g=1\), et
\begin{equation}
 |f-f_n|=|g-g_n|\leq\frac{C}{M_0(B-1)B^n}.
 \label{x_reciprocidad:lib03:cota-cola}
\end{equation}
\end{theorem}

\begin{proof}
La queue absolue satisfait
\[
 \sum_{j>n}\frac{|c_j|}{B^j}
 \leq C\sum_{j>n}B^{-j}=\frac{C}{(B-1)B^n}.
\]
La convergence de la série géométrique et
\eqref{x_reciprocidad:lib03:unidad-normalizada} prouvent toutes les affirmations.
\end{proof}

Pour des retenues des deux signes, une condition suffisante de positivité
à toutes les étapes est \(x_0,y_0\geq C/(B-1)\). Pour l’alphabet
canonique, \(x_0\geq1,y_0\geq0\) suffit. Dans ce second cas,
\(0\leq\theta\leq1\), \(f_n\) décroît et \(g_n\) croît, et
\begin{equation}
 0\leq f_n-f=g-g_n\leq\frac1{M_0B^n}.
 \label{x_reciprocidad:lib03:cota-canonica}
\end{equation}
La somme demeure fixe et les proportions individuelles évoluent.
L’histoire finie reste récupérable lorsque l’on conserve les
coordonnées entières et l’échelle.

\section{Cylindres compatibles et multiplicité de frontière}
\label{x_reciprocidad:lib03:cilindros}

Considérons l'alphabet canonique complet et une entrée \(x_0\geq1,y_0\geq0\). Un préfixe de profondeur \(n\), de cumul \(C_n\), admet exactement les valeurs limites
\begin{equation}
 I_n(C_n)=\left[\frac{C_n}{B^n},\frac{C_n+1}{B^n}\right]
 \label{x_reciprocidad:lib03:cilindro}
\end{equation}
de la lecture \(\theta\).

\begin{theorem}[Cylindres emboîtés et distinction des histoires]
\label{x_reciprocidad:lib03:teo-cilindros}
Les \(B\) cylindres enfants de \(I_n(C_n)\) recouvrent le parent et ne se
chevauchent qu’aux extrémités. Les cylindres d’une histoire infinie
ont pour diamètre \(B^{-n}\) et une intersection réduite à un seul point.
Deux histoires canoniques différentes produisent la même valeur si et seulement
si, à leur première position différente, leurs chiffres diffèrent de un et les
queues ultérieures sont, respectivement, entièrement \(B-1\) pour le chiffre
inférieur et entièrement nulles pour le supérieur.
\end{theorem}

\begin{proof}
L'enfant associé à \(c\) a pour extrémités \((BC_n+c)/B^{n+1}\) et \((BC_n+c+1)/B^{n+1}\). Ces expressions démontrent l'inclusion, la couverture et le recouvrement uniquement aux extrémités. La queue d'un développement parcourt l'intervalle \([0,B^{-n}]\) ; le choix successif d'un des enfants produit un développement de chaque point de cet intervalle. Les intervalles fermés emboîtés ont des diamètres qui tendent vers zéro, ce qui prouve l'existence et l'unicité de l'intersection pour chaque histoire.

Soient \(k\) la première position distincte et \(d_k>c_k\). Si les
lectures coïncident,
\[
 \frac{d_k-c_k}{B^k}
 =\sum_{j>k}\frac{c_j-d_j}{B^j}
 \leq\sum_{j>k}\frac{B-1}{B^j}=\frac1{B^k}.
\]
La différence positive entière \(d_k-c_k\) doit être un. L’égalité
dans la borne de la queue exige \(c_j=B-1\), \(d_j=0\) pour tout
\(j>k\). Réciproquement, ces queues ont pour différence exactement
\(B^{-k}\), ce qui compense le premier chiffre. La
caractérisation complète est prouvée.
\end{proof}

L’image du cylindre pour \(f\) a pour diamètre
\(1/(M_0B^n)\), qui coïncide avec \eqref{x_reciprocidad:lib03:cota-canonica}.
Les histoires \((1,0,0,\ldots)\) et
\((0,B-1,B-1,\ldots)\) produisent toutes deux \(\theta=1/B\), en conservant
des préfixes et des états entiers finis distincts. Une marque de frontière
maintient les deux provenances. Une convention qui exclut les queues ultimement
constantes \(B-1\) sélectionne un développement pour \(\theta<1\) ;
l’extrémité \(\theta=1\) est conservée séparément, car dans
l’alphabet fractionnaire, elle correspond à la queue constante \(B-1\).

L'histoire des opérations, le système de cylindres et la valeur limite sont des objets reliés au moyen de ces applications. L'unicité de la lecture d'une histoire coexiste avec la multiplicité des histoires aux valeurs de frontière.

\section{Retenues signées et spécification prospective}
\label{x_reciprocidad:lib03:firmados}

Permettre des retenues hors de l’alphabet canonique exige de conserver leurs
relèvements. Dès deux étapes, les mots
\[
 (c_1,c_2)=(1,-1),\qquad(d_1,d_2)=(0,B-1)
\]
produisent \(C_2=B-1\) et la même paire finale. Tous deux sont bornés par
\(B-1\) et peuvent préserver la positivité, par exemple à partir de
\((73,27)\) en base dix. Le mot signé, ou les relèvements de
\eqref{x_reciprocidad:lib03:elevacion}, distingue cette information. Le théorème de
convergence reste valable pour les deux mots ; l’inverse canonique
possède, en revanche, le domaine indiqué dans son énoncé.

La transition \eqref{x_reciprocidad:lib03:caso-729} fixe \(c_1=1\). Une règle supplémentaire qui impose \(c_j=1\) pour tout \(j\) produit
\begin{equation}
 \theta_n=\frac{1-B^{-n}}{B-1},\qquad
 \theta=\frac1{B-1}.
 \label{x_reciprocidad:lib03:acarreo-constante}
\end{equation}
En base dix, à partir de \((73,27)\), la suite est
\[
 (73,27)\longmapsto(729,271)\longmapsto(7289,2711)
 \longmapsto(72889,27111),
\]
et la première fraction converge vers
\begin{equation}
 \frac{73}{100}-\frac1{900}=\frac{164}{225}
 =0{,}728888\ldots.
 \label{x_reciprocidad:lib03:limite-constante}
\end{equation}
La conclusion appartient à cette règle supplémentaire de retenues constantes. Avec seulement \(c_1=1\) et des prolongements canoniques non négatifs, le cylindre de la première étape fournit exactement
\begin{equation}
 \frac{b^3-1}{B^3}\leq f\leq\frac{b^3}{B^3}.
 \label{x_reciprocidad:lib03:primer-cilindro}
\end{equation}
En base dix, c’est \([0{,}728,0{,}729]\). Changer le signe ou l’alphabet
change cette famille et requiert de conserver la mémoire supplémentaire
correspondante. La présence de \(729\) dans la complétion détermine
la transition arithmétique écrite ; une identification à un autre lecteur de
constantes utilise en outre sa règle effective de génération et son
entrelacement.

\section{Réalisation opératorielle et amplitudes de transport}
\label{x_reciprocidad:lib03:operativo}

La procédure résultante comprend la mise à jour au moyen de \eqref{x_reciprocidad:lib03:lc}, la vérification locale par résidu, la reconstruction de tout mot canonique au moyen de \eqref{x_reciprocidad:lib03:recuperacion-digitos}, le contrôle de la lecture limite au moyen de \eqref{x_reciprocidad:lib03:cota-cola} et le regroupement d'événements au moyen de \eqref{x_reciprocidad:lib03:ley-bloques}. Chaque opération maintient l'ordre et le compteur d'échelle qu'utilise son inverse.

La conservation de la somme et la conservation quadratique ont des formes
différentes. Pour un transfert normalisé
\(\delta=c/(BM)\),
\begin{equation}
 (f-\delta)^2+(g+\delta)^2-(f^2+g^2)
 =2\delta(g-f)+2\delta^2.
 \label{x_reciprocidad:lib03:variacion-cuadratica}
\end{equation}
Cette expression détermine la variation quadratique et est généralement
différente de zéro. Le lien avec une isométrie s’obtiendra sur les
amplitudes des étiquettes admissibles, au moyen de la bijection par
résidus de la \cref{x_reciprocidad:lib05:levantamiento}. Ainsi, le transport arithmétique
et sa réalisation unitaire se composent en conservant leurs quantités
respectives, au lieu de les identifier par une coïncidence scalaire.
